\PageTitle{03 / Benchmark observations}{Overall read latency}
All 100 raw samples have \texttt{ok=true} and \texttt{error=null}. The supplied values reproduce under nearest-rank calculation. Figure~\ref{fig:cdf} uses the complete ordered sample, rather than a fitted distribution.

\vspace{3mm}
\Metric{100/100}{Completed successfully}{No recorded tool errors}\hfill
\Metric{11.946 s}{Total wall time}{Includes pacing}\hfill
\Metric{8.37/s}{Completion rate}{Not a capacity estimate}
\vspace{4mm}

\begin{figure}[h]\centering\begin{tikzpicture}
\begin{axis}[xmin=0,xmax=750,ymin=0,ymax=100,xlabel={Client-observed latency (ms)},ylabel={Calls at or below latency (\%)},grid=major,ytick={0,25,50,75,100},height=72mm]
\addplot[ConversationBlue,very thick,const plot] coordinates {(0,0) (205,1.0) (214,2.0) (217,3.0) (219,4.0) (222,5.0) (226,6.0) (229,7.000000000000001) (232,8.0) (239,9.0) (240,10.0) (243,11.0) (245,12.0) (246,13.0) (251,14.000000000000002) (251,15.0) (253,16.0) (256,17.0) (264,18.0) (265,19.0) (269,20.0) (269,21.0) (273,22.0) (273,23.0) (275,24.0) (276,25.0) (278,26.0) (278,27.0) (281,28.000000000000004) (287,28.999999999999996) (292,30.0) (302,31.0) (310,32.0) (316,33.0) (324,34.0) (328,35.0) (330,36.0) (334,37.0) (337,38.0) (340,39.0) (342,40.0) (345,41.0) (347,42.0) (349,43.0) (351,44.0) (352,45.0) (353,46.0) (354,47.0) (357,48.0) (358,49.0) (361,50.0) (363,51.0) (366,52.0) (369,53.0) (370,54.0) (371,55.00000000000001) (372,56.00000000000001) (372,56.99999999999999) (373,57.99999999999999) (377,59.0) (377,60.0) (380,61.0) (383,62.0) (384,63.0) (384,64.0) (385,65.0) (390,66.0) (391,67.0) (393,68.0) (396,69.0) (397,70.0) (401,71.0) (402,72.0) (406,73.0) (407,74.0) (407,75.0) (411,76.0) (411,77.0) (416,78.0) (450,79.0) (465,80.0) (473,81.0) (489,82.0) (494,83.0) (500,84.0) (501,85.0) (506,86.0) (513,87.0) (518,88.0) (611,89.0) (619,90.0) (628,91.0) (634,92.0) (643,93.0) (656,94.0) (662,95.0) (668,96.0) (675,97.0) (687,98.0) (691,99.0) (695,100.0)};
\end{axis}
\end{tikzpicture}
\caption{Empirical cumulative distribution of 100 client-observed read durations. The x-axis starts at zero; the shortest observed call was 205 ms.}\label{fig:cdf}\end{figure}
\Source{S2; recomputed by \texttt{analyze\_evidence.py}}

\begin{table}[h]\caption{Overall latency summary}\EvidenceTableStart
\begin{tabularx}{\linewidth}{@{}Xrrrr@{}}\toprule
Metric & Median & p95 & p99 & Maximum\\\midrule
Client latency, ms & 361 & 662 & 691 & 695\\
Ordered rank / 100 & 50 & 95 & 99 & 100\\\bottomrule
\end{tabularx}\end{table}

\subsection*{Interpretation}
The run demonstrates completion of this short, small-result read workload. Tail latency is materially above the median within the sample. No latency service-level objective was supplied, so these values are observations rather than a pass/fail against a service target.

\Limit{A 100/100 observed success count is not a population reliability estimate or production guarantee. There is no direct-HTTP control in this run. These values cannot isolate bridge overhead, server processing time, network delay, or database time.}
